Safety evaluations report nulls constantly: ``we could not elicit the behaviour.'' That
sentence is load-bearing in deployment decisions, and it is ambiguous, because it can mean the
model is robust or that the experiment was badly designed. A popular remedy is adaptivity: treat
a null as feedback and redesign. We test this directly, holding total compute constant, which no
prior work does. We replay four protocol families against 437k real per-iteration attack traces,
a complete $5\times4\times100$ cost-to-discovery matrix, and a live grid on three open-weights
instruct models under an anti-sycophancy intervention. Adaptivity helps enormously: a protocol
that spreads a fixed budget over three design families detects 100\% of behaviours where the
default single-design protocol detects 0\% at every budget up to 10{,}000 iterations. But almost
none of that gain comes from the proposed mechanism. Against a comparator that switches design
family on a blind fixed schedule with identical budget and identical designs, using the null as
a trigger gains nothing ($\Delta = +0.02$) and loses at large budgets ($\Delta = -0.12$). We
explain why: conditional on a run still being null, the continuous progress signal predicts
eventual success at AUC $0.45$--$0.61$, with every interval covering chance. Adaptivity at the
portfolio level does work ($+0.65$ behaviours on GPT-4; $+0.72$ live on Qwen2.5-7B). Finally,
nominally 95\% bounds on residual vulnerability achieve 63--67\% coverage, and an anytime-valid
correction barely helps, because the estimand, not the stopping rule, is wrong.
